\chapter{SMPL--Anny joint correspondence}
\label{app:joints}
% Never \input tables/draft_mapping.tex: it defines the same table label as this file.
% TODO (author): check the spine rows, neck offset and scale factor against a converter log (README.md).

The converter has to know which Anny bone stands for each SMPL joint. The two skeletons are the ones in
Table~\ref{tab:body-models}: SMPL has \BgSmplJoints{} joints and Anny has \BgAnnyBones{} bones, of which the root and
\DataFittedBones{} rotate in the fit. Some joints of the two rigs sit at different places, so names alone are not
enough. The converter compares the two rest skeletons (the default standing pose) before it fits the posed frames.

For each hip, Anny has two possible bones, the upper-leg bone and the pelvis bone. I take the one whose head (the
point where the bone starts) lies about as far from the pelvis as SMPL's hip joint; for the v23 fits this is the
upper-leg bone. For the spine, I order the
Anny spine bones by their height, scaled so that the hip midpoint is 0 and the neck is 1. I then keep the ordered
triple of bones whose heights are together closest to the heights of SMPL's three spine joints. For the limbs and the
head, the converter takes the Anny bone of the same body part from a fixed list of names: the shoulder, for example,
goes to the head of the upper-arm bone. Each run saves the bones it picked together with the fits.

For the torso joints, the converter fits a virtual joint (Equation~\eqref{eq:virtual-joint}): a point that sits at a
fixed offset from the Anny bone head. The converter measures this offset once, in the rest pose. Let $\mathbf j^S_k$
be the position of SMPL joint $k$, $\mathbf o_A$ and $\mathbf o_S$ the hip midpoints of the Anny and the SMPL rest
skeleton, $s$ the ratio of Anny's to SMPL's hip-midpoint-to-neck length, $\mathbf Q_{S\to A}$ the matrix that lines
up the axes of the two rest skeletons, and $\mathbf h^0_b$ and $\mathbf Q^0_b$ the rest head and orientation of the
bone $b$ assigned to joint $k$:
\begin{equation}
 \mathbf j^{S\to A}_k=\mathbf o_A+s\mathbf Q_{S\to A}(\mathbf j^S_k-\mathbf o_S),\qquad
 \boldsymbol\delta_k=(\mathbf Q^0_b)^\top(\mathbf j^{S\to A}_k-\mathbf h^0_b).
 \label{eq:rest-offset}
\end{equation}
In words, the first formula moves the SMPL joint into Anny's rest skeleton. It shifts the joint from SMPL's hip
midpoint to Anny's, scales it by $s$ and turns it with $\mathbf Q_{S\to A}$. The second formula then measures the
offset from the Anny bone head, in that bone's own frame. The offset stays fixed and
turns with the bone, so it is not an extra joint fitted in each frame.

\begin{table}[htbp]
\centering\small
\begin{tabular}{@{}>{\raggedright}p{3.4cm}>{\raggedright\arraybackslash}p{8.7cm}@{}}
\toprule
SMPL target & Anny correspondence and fitting point \\
\midrule
Pelvis & root bone with a measured torso offset \\
Left/right hip & upper-leg or pelvis bone, chosen by its distance from the pelvis (the upper-leg bone in v23),
with an offset \\
Spine1, spine2, spine3 & spine bones matched by normalised rest height, in order, with offsets \\
Neck & neck candidate with a measured rest-frame offset \\
Left/right collar & clavicle bone with an offset \\
Shoulder, elbow, wrist & heads of the matching upper-arm, lower-arm and wrist bones \\
Knee, ankle, foot & heads of the matching lower-leg, foot and toe bones \\
Head and hand targets & heads of the matching bones; their pose follows from the parent bones that the fit is
allowed to rotate \\
\bottomrule
\end{tabular}
\caption[Construction of the SMPL--Anny correspondence]{Correspondence used by the converter for the v23 labels.
The exact bone names and offsets are saved with each fit and can differ between Anny rig versions. The fit rotates the
root and \DataFittedBones{} body bones; the others, including most fingers and the finer head and neck bones, keep
their rest rotation.}
\label{tab:mapping}
\end{table}

Each fit stores its own correspondence in three arrays: \path{anny_bone_names}, \path{joint_offsets_local} and
\path{offset_joints}. To reproduce a particular fit, use these arrays and not this table. The code is in
\path{scripts/data/smpl_to_anny.py}. Appendix~\ref{app:conversion} shows how geometric matching and virtual joints
changed the fit error, and why matching by name failed in my first converter.
